\documentclass{article}
\begin{document}
\title{Forcing one sensor's capture is not the same invocation as forcing the whole manager}
\maketitle

\paragraph{Scope.}
Run the selected pinned Sensor::Update, Manager::RunOnce and
Manager::Sensor method bodies. The fixture supplies three loaded
sensors, exact time and callback accounting. Sensor work weights
and deadlines are model inputs, not measured render costs.
The public C++ getter gives the local caller one sensor pointer.
No external permission or actual deployed controller is assumed.

\paragraph{Models and comparisons.}
The target must furnish one capture. The other two sensors'
regularly due tasks must still run at this cut. Their not-yet-due
captures are optional. An ordinary manager call invokes the
native gate for each member and can fail to provide the target
when it is not due. The native manager's force flag goes to
every member, not just the target; that is a different workload.
Its forced callbacks do not advance nominal scheduling.
A scoped implementation runs the ordinary cut and forces the
target only if its qualified result is still missing.
A negative control that forces first and then runs ordinary
processing can generate two target captures when it was due.

When callback cost is additive, the minimum work to return
with the named result set is the sum over the target and the
other due sensors. Scoped conditional forcing attains that
cost for every admitted case. It is not an advance over an
equally informed direct gate. If the receiver can act on
early callbacks instead of waiting for manager return, an
optional extra update after all required results does not
cause a deadline miss at that cut, though it still costs work.
If it occupies the earlier global schedule it can matter.
The ordering of the actual loaded IDs is therefore recorded.

\paragraph{Results and residuals.}
Test all three-member orders, due masks and costs 1/2. In
1,536 cases, wide force exceeds the minimal return-bound
budget in 288; 96 of those can still supply all named
results by budget under an early-callback law. The unconditional
target-first control duplicates target output in 192 cases.
The scoped control keeps all due work and has no duplicates.
The model assumes synchronous access to the member callback
result at the accounting cut, and no safety assertion follows
for other consumers of unsolicited extra frames. No actual
camera acquisition, force availability under the stock mission,
or independent architecture acceptance is asserted.
\end{document}
